% Generated draft: requires mathematical and source review.
\documentclass[11pt]{article}
\usepackage[margin=25mm]{geometry}
\usepackage{fontspec}
\setmainfont{Latin Modern Roman}
\usepackage{amsmath,amssymb,mathtools,unicode-math}
\setmathfont{Latin Modern Math}
\usepackage{enumitem,needspace,xurl,hyperref}
\setlength{\parindent}{0pt}
\setlength{\parskip}{5pt}
\setlength{\emergencystretch}{3em}
\newcommand{\N}{\mathbb{N}}
\newcommand{\Z}{\mathbb{Z}}
\newcommand{\Q}{\mathbb{Q}}
\newcommand{\R}{\mathbb{R}}
\newcommand{\C}{\mathbb{C}}
\newcommand{\F}{\mathbb{F}}
\newcommand{\A}{\mathbb{A}}
\newcommand{\PP}{\mathbb{P}}
\newcommand{\E}{\mathbb{E}}
\newcommand{\eps}{\varepsilon}
\newcommand{\dd}{\,\mathrm{d}}
\begin{document}
\section*{UnsolvedMath ID: 700311}
\textbf{Source problem code:} catalogue entry
\subsection{Definitions and mathematical statement}
Let \(p\) be a prime. A pro-\(p\) group is a topological group isomorphic to an inverse limit of finite \(p\)-groups. Such a group \(G\) is torsion-free if its identity is its only element of finite order. Define its continuous \(p\)-cohomological dimension by
\[
\operatorname{cd}_p(G)=\sup\bigl\{n\geq 0: H^n_{\mathrm{cont}}(G,M)\neq 0\text{ for some discrete }p\text{-primary torsion }G\text{-module }M\text{ with continuous action}\bigr\}.
\]
The source writes \(\operatorname{cd}G\); here this is interpreted as \(\operatorname{cd}_p(G)\).

A pro-\(p\) presentation with one defining relator has the form
\[
G\cong F/R,\qquad R=\overline{\langle\!\langle r\rangle\!\rangle}_F,
\]
where \(F\) is a free pro-\(p\) group and \(R\) is the smallest closed normal subgroup of \(F\) containing \(r\). To express an essential single relation precisely, one may require \(R\neq\{1\}\) and require the presentation to be minimal: the induced map
\[
F/\Phi(F)\longrightarrow G/\Phi(G)
\]
is an isomorphism, where \(\Phi(H)=\overline{H^p[H,H]}\) is the Frattini subgroup of a pro-\(p\) group \(H\).

\textbf{Requested assertion, under this essential-relation interpretation.} For every prime \(p\) and every torsion-free pro-\(p\) group \(G\) having such a minimal presentation with a nontrivial relation kernel normally generated by one element,
\[
\operatorname{cd}_p(G)=2.
\]
\textbf{Unresolved source qualifications.} The supplied statement says only ``a single defining relation'' and does not explicitly impose minimality or nonredundancy. The precise presentation convention above is therefore an interpretation requiring review, rather than an additional hypothesis explicitly supplied by the source. No finite-generation hypothesis is supplied, and none is imposed here. The source also reports that the assertion was proved for a large class of groups in the cited 1967 paper of J. P. Labute, without specifying that class.
\subsection{Short English statement}
H. Koch asks whether every torsion-free pro-\(p\) group with one essential defining relation has continuous \(p\)-cohomological dimension exactly two. The intended presentation convention needs confirmation. The supplied registry labels the problem open and reports a partial result of Labute; these are supplied source assessments, not independently verified conclusions.
\subsection{Sources}
\begin{itemize}
\item Local UnsolvedMath registry, permanent ID \(700311\), source code \texttt{catalogue entry}, title \emph{Kourovka Notebook Problem 5.27}. The supplied record labels the problem \texttt{open}; its update date is \texttt{2026-10-16T14:00:00Z}. This status is supplied metadata and has not been independently verified.
\item \emph{The Kourovka Notebook (Unsolved Problems in Group Theory)}, original problem 5.27, fifth issue (1976), proposed by H. Koch. The supplied background identifies the consulted edition as the twenty-first issue edition (2026), \texttt{arXiv:1401.0300v48}. The supplied collection description names E. I. Khukhro and V. D. Mazurov as editors. These bibliographic details are reproduced from the record without independent source consultation.
\item J. P. Labute, \emph{Inv. Math.}, volume 4, number 2 (1967), pp. 142--158. The supplied statement attributes a proof for a large class of groups to this reference; it supplies neither an article title nor the hypotheses describing that class.
\end{itemize}
\end{document}
